\documentclass[11pt]{article}

% horizontal margins: 1.0 + 6.5 + 1.0 = 8.5
\setlength{\oddsidemargin}{0.0in} \setlength{\textwidth}{6.5in}
% vertical margins: 1.0 + 9.0 + 1.0 = 11.0
\setlength{\topmargin}{-0.35in} \setlength{\headheight}{12pt}
\setlength{\headsep}{13pt} \setlength{\textheight}{675pt}
\setlength{\footskip}{24pt}
\thispagestyle{empty}


\usepackage{marvosym}
\usepackage[plainpages=false, colorlinks=true,
   citecolor=blue, filecolor=blue, linkcolor=blue,
   urlcolor=blue]{hyperref}

% Section titles with a thin rule underneath
\usepackage{titlesec}
\titleformat{\section}{\Large\bfseries}{}{0pt}{}[\vspace{2pt}\titlerule]
\titlespacing*{\section}{0pt}{12pt}{6pt}

\begin{document}
\emergencystretch=2em % let LaTeX stretch spacing slightly instead of overflowing long titles

%#################################################################################
% HEADER
\begin{center}
  {\bfseries \Large Irabiel Romero} \\[0.1in]
  % Applied Mathematician \\
  % Ph.D. Candidate in Applied Mathematics \\
  % University of California, Merced \\
  \href{mailto:irabielr96@gmail.com}{irabielr96@gmail.com} \quad
  (661) 240-7760 \quad
  \href{https://www.linkedin.com/in/irabiel}{linkedin.com/in/irabiel} \quad
  \href{https://irabiel.github.io}{irabiel.github.io}
\end{center}

\vskip 0.2in

%#################################################################################
\section*{Research Interests}

\noindent
\begin{tabular}{lp{0.75\textwidth}}
General: & Mathematical modeling, numerical analysis, large-scale
inverse problems, PDE-constrained optimization, Bayesian
inference, uncertainty
quantification, optimal experimental design, optimization under
uncertainty, numerical optimization, numerical linear algebra,
multigrid solvers, scientific machine learning \\
Applications: & Power-grid modeling
\end{tabular}

%#################################################################################
\section*{Education}

\vskip 0.1in
\noindent
\begin{tabular}{lp{0.82\textwidth}}
  Dec. 2026 & {\bf Ph.D. Applied Mathematics}, University of California,
  Merced, CA, USA (expected). Thesis: \emph{Scalable Methods
  for Estimation of Power Systems under Uncertainty} (Advisors: Prof.
  Noemi Petra and Prof. Roummel Marcia) \\
  May 2021 & {\bf M.S. Applied Mathematics}, California State
  University, Long Beach, CA, USA. Thesis: \emph{Numerical Methods for
  Multiscale Modeling in Bacterial Colony Growth} (Advisor: Prof. Hui Sun) \\
  May 2018 & {\bf B.S. Applied Mathematics} and {\bf B.A. Physics},
  California State University, Long Beach, CA, USA \\
  May 2016 & {\bf A.S. Mathematics}, Bakersfield College, Bakersfield,
  CA, USA
\end{tabular}

%#################################################################################
\section*{Ph.D.'s Research Experience}

\begin{itemize}
\item {\bf Graduate Research Intern}, Center for Applied Scientific
  Computing (CASC), Lawrence Livermore National Laboratory, Livermore, CA,
  June 2026--Aug. 2026
  \begin{itemize}
  \item Integrated a machine learning surrogate into the transient
    stability-constrained AC optimal power flow (TS-ACOPF) solver,
    replacing stability constraints that are computationally intractable
    to enforce directly.
  \item Designed and ran validation tests comparing surrogate-constrained
    solutions against full simulation.
  \end{itemize}
\item {\bf Graduate Student Researcher}, Lawrence Livermore National
  Laboratory, Livermore, CA, Spring 2025--Spring 2026 (Project
  manager: Dr. Cosmin Petra)
  \begin{itemize}
  \item Extended LLNL's security-constrained AC optimal power flow
    (ACOPF) solver, ExaJuGO, to analyze wildfire-driven N-k
    contingencies on the California Test System, and developed
    interactive visualization tools for regional and system-wide
    analysis.
  \item Developed a convolutional neural network (CNN) surrogate for
    generator stability and integrated it into ACOPF to enforce
    transient stability constraints that would otherwise be
    computationally intractable.
  \end{itemize}
\item {\bf Data Science Summer Institute}, Lawrence Livermore National
  Laboratory, Livermore, CA, May 2024--Aug. 2024 (Project manager:
  Dr. Ignacio Aravena)
  \begin{itemize}
  \item Developed a general parallel framework (C++, MPI, HDF5) using
    Particle Swarm Optimization to generate DC power flow databases for
    studying the effects of line susceptance on power flows and voltage
    angles.
  \item Developed CNN models to classify heartbeats from
    electrocardiography (ECG) signals, identify the type of
    irregularity, and estimate cardiac activation times.
  \end{itemize}
\item {\bf Data Science Challenge Internship}, Lawrence Livermore
  National Laboratory, Livermore, CA, July 2023--Aug. 2023 (Project
  manager: Dr. Mikel Landajuela; project support: Brian Gallagher and
  Cindy Gonzales)
  \begin{itemize}
  \item Served as a team lead in developing a 1D CNN that reconstructs
    the heart's transmembrane voltages and activation times from ECG
    signals.
  \end{itemize}
\item {\bf Graduate Student Researcher}, University of California,
  Merced, CA, Jan. 2023--current (Advisors: Prof. Noemi Petra and
  Prof. Roummel Marcia)
  \begin{itemize}
  \item Developed a Bayesian inverse problem governed by the DC power
    flow equations to estimate voltage angles and uncertain line
    parameters, and an optimal experimental design method for PMU
    placement.
  \item Led a project with Dr.\ Cosmin G.\ Petra (CASC) and Dr.\ Kathleen
    Schmidt (Uncertainty Quantification group lead), LLNL, Fall 2026:
    developed an ML surrogate for TS-ACOPF that learns the distribution of
    transient stability index (TSI) values over all possible bus fault
    locations and uses it to estimate the probability that the system
    becomes unstable; added conformal prediction to give calibrated,
    distribution-free bounds on that probability.
  \end{itemize}
\item {\bf Graduate Student Researcher}, University of California,
  Merced, CA, May 2022--Aug. 2022 (Advisors: Prof. Noemi Petra and
  Prof. Maxime Theillard)
  \begin{itemize}
  \item Developed a generic algebraic multigrid solver for M-matrices,
    and an inverse problem governed by an advection--diffusion equation
    to recover the diffusion parameter from sparse data.
  \end{itemize}
\end{itemize}

%#################################################################################
\section*{Master's Research Experience}

\begin{itemize}
\item {\bf Graduate Researcher}, California State University, Long
  Beach, CA, May 2020--May 2021 (Advisor: Prof. Hui Sun)
  \begin{itemize}
  \item Extended the modified Newton--multigrid solver to a 2D
    reaction--diffusion model of bacterial crossfeeding, and coupled it
    with a kinetic Monte Carlo model for agent-based cell simulation.
  \end{itemize}
\item {\bf Graduate Researcher}, California State University, Long
  Beach, CA, June 2019--Aug. 2019 (Advisor: Prof. Hui Sun)
  \begin{itemize}
  \item Developed a modified Newton--multigrid solver for a stiff system
    of nonlinear reaction--diffusion partial differential equations
    (PDEs) describing chemical concentrations in a bacterial colony.
  \end{itemize}
\item {\bf Graduate Researcher}, University of California, San Diego,
  CA, June 2018--Aug. 2018 (Advisor: Prof. Hui Sun)
  \begin{itemize}
  \item Extended a 3D computational bacterial colony model with PDEs for
    additional chemical species that influence cell metabolism,
    movement, and colony morphology.
  \end{itemize}
\end{itemize}

%#################################################################################
\section*{Publications}

\begin{itemize}
\item {\bf I. Romero}, N.-Y. Chiang, I. Aravena, R. Marcia, N. Petra,
  and C.~G. Petra. ``Real-time power system machine learning-assisted
  transient stability preventive control''. In preparation.
\item {\bf I. Romero}, H. Kannan, A. Valdez, M. Winger, B. Li, and
  H. Sun. ``Advanced numerical methods for reaction-diffusion equations
  modeling the growth of bacterial colonies''. SSRN preprint (2026)
\item {\bf I. Romero}, N.-Y. Chiang, I. Aravena, R. Marcia, N. Petra,
  J.-P. Watson, and C.~G. Petra. ``Wildfire-driven N-k contingency
  analysis for security-constrained AC optimal power flow in the
  California power grid''. IEEE Electrical Power and Energy Conference
  (EPEC) (2026)
\item {\bf I. Romero}, R. Marcia, I. Aravena, and N. Petra. ``Optimal
  PMU placement for state estimation with grid parameter uncertainty''.
  2025 IEEE PES Grid Edge Technologies Conference and Exposition (Grid
  Edge), IEEE (2025)
\item A. Ranganath, {\bf I. Romero}, M. Singhal, and R. Marcia.
  ``Quasi-Adam: Accelerating Adam using quasi-Newton approximations''.
  2024 International Conference on Machine Learning and Applications
  (ICMLA), IEEE (2024)
\item {\bf I. Romero}. ``Numerical methods for multiscale
  modeling in bacterial colony growth''. Master's Thesis, California State University, Long
  Beach (2021)
\end{itemize}

%#################################################################################
\section*{Research Support}

\vskip 0.1in
\noindent
\begin{tabular}{rp{0.8\textwidth}}
  2025 & LLNL subcontract, Lawrence Livermore National Laboratory (PI:
  Dr. Cosmin G. Petra; with Dr. Nai-Yuan Chiang), Spring--Summer 2025 \\\\
  2023--2026 & Graduate Student Researcher, NSF CAREER \#1654311 and NSF
  AMPS \#229495 \\\\
  2022 & Summer Research, Fields Family Fellowship and Applied Math
  Graduate Program Summer Research Fellowship, Summer 2022
\end{tabular}

%#################################################################################
\section*{Teaching Experience}

\vskip 0.15in
\noindent
\begin{tabular}{llp{0.58\textwidth}}
  Fall 2022 & {\bf Teaching Assistant} & {\bf Math 131} - {\it Numerical
    Methods for Scientists and Engineers}, University of California,
  Merced\\
  Spring 2022 & {\bf Teaching Assistant} & {\bf Math 12} - {\it Calculus
    II}, University of California, Merced\\
  Fall 2021 & {\bf Teaching Assistant} & {\bf Math 24} - {\it Linear
    Algebra and Differential Equations}, University of California,
  Merced\\
  Fall 2020 & {\bf Teaching Assistant} & {\bf Math 112} - {\it Essential
    Algebra A}, California State University, Long Beach
\end{tabular}

%#################################################################################
\section*{Oral Presentations}

\begin{itemize}
\item {\it Scalable Methods for Estimation of Power Systems under
  Uncertainty}, Optimization Seminar, University of California, Merced,
  CA, Fall 2026
\item {\it Real-Time Power System Machine Learning-Assisted Transient
  Stability Preventive Control}, Optimization Seminar, University of
  California, Merced, CA, Spring 2026
\item {\it Wildfire-driven N-k contingency analysis for
  security-constrained AC optimal power flow in the California power
  grid}, Optimization Seminar, University of California, Merced, CA,
  Fall 2025
\item {\it Optimal PMU Placement for State Estimation with Grid
  Parameter Uncertainty}, Optimization Seminar, University of
  California, Merced, CA, Spring 2025
\item {\it Optimal PMU Placement for State Estimation with Grid
  Parameter Uncertainty}, 28th Annual CASIS Workshop, Livermore, CA,
  Summer 2024
\item {\it Embedding Power Flow into Machine Learning for Parameter and
  State Estimation}, Optimization Seminar, University of California,
  Merced, CA, Fall 2022
\item {\it Numerical Methods for Multiscale Modeling in Bacterial Colony
  Growth}, 32nd Annual CSULB Student Research Competition, Long Beach,
  CA, Spring 2020
\item {\it Spatiotemporal Establishment of Growing Bacterial Colonies},
  31st Annual CSULB Student Research Competition, Long Beach, CA, Spring
  2019
\item {\it Spatiotemporal Establishment of Growing Bacterial Colonies},
  UCSD Q-Bio Seminar, San Diego, CA, Summer 2018
\end{itemize}

%#################################################################################
\section*{Poster Presentations}

\begin{itemize}
\item {\it Real-Time Power System Machine Learning-Assisted Transient
  Stability Preventive Control}, 3rd SIAM Northern and Central
  California Sectional Conference ({\bf NCC26}), Davis, CA, Fall 2026
  (Poster Award)
\item {\it Wildfire-driven N-k contingency analysis for
  security-constrained AC optimal power flow in the California power
  grid}, 2nd SIAM Northern and Central California Sectional Conference
  ({\bf NCC25}), Berkeley, CA, Fall 2025 (Poster Award)
\item {\it Optimal PMU Placement for State Estimation with Grid
  Parameter Uncertainty}, IEEE PES Grid Edge Technologies Conference and
  Exposition, San Diego, CA, Spring 2025
\item {\it Optimal PMU Placement for State Estimation with Grid
  Parameter Uncertainty}, 28th Annual CASIS Workshop, Livermore, CA,
  Summer 2024
\item {\it Hearts with Machine Learning}, LLNL Data Science Challenge,
  Livermore, CA, Summer 2023
\item {\it Numerical Methods for Multiscale Modeling in Bacterial Colony
  Growth}, Graduate Research Conference, Long Beach, CA, Fall 2019
\item {\it Spatiotemporal Establishment of Growing Bacterial Colonies},
  Graduate Research Conference, Long Beach, CA, Fall 2018
\end{itemize}

%#################################################################################
\section*{Selected Honors and Awards}

\noindent
\begin{tabular}{@{}l@{\hspace{1em}}l}
  2026 & Outstanding Poster Award, SIAM NCC26, Davis, CA \\
  2025 & Best Poster Award, SIAM NCC25, Berkeley, CA \\
  2021 & Fields Family Fellowship, UC Merced \\
  2021 & Outstanding Thesis Award, CSU Long Beach \\
  2020 & Richard Green Graduate Fellowship, CSU Long Beach \\
  2019 & Summer Student Research Assistantship, CSU Long Beach \\
  2019 & Samuel Councilman Scholarship, CSU Long Beach \\
  2018 & President's List, CSU Long Beach \\
  2017 & Dean's List, CSU Long Beach
\end{tabular}

%#################################################################################
\section*{Professional Development}

\begin{itemize}
\item {\it Uncertainty Quantification Summer School}, University of
  Southern California, Los Angeles, CA, August 2026
\item {\it Structure-Preserving Scientific Computing and Machine
  Learning Summer School and Hackathon} (PIMS CRG), University of
  Washington, Seattle, WA, June 2025
\item {\it FEniCS 2022 Conference}, University of California, San
  Diego, La Jolla, CA, August 2022
\end{itemize}

%#################################################################################
\section*{Graduate Coursework}

\noindent
\begin{tabular}{lp{0.75\textwidth}}
UC Merced: & Numerical Optimization (Fall 2023), Scientific Computing
(Fall 2022), Advanced Methods of Applied Mathematics (Spring 2022),
Numerical Solution of Differential Equations I and II (Fall 2021,
Spring 2022), Partial Differential Equations (Spring 2022),
Computational Inverse Problems (Fall 2021), Advanced Ordinary
Differential Equations (Fall 2021) \\
CSU Long Beach: & Finite Element Method (Spring 2021), Mathematical
Modeling with Graph Theory (Spring 2020), Sobolev Spaces (Fall 2019),
Partial Differential Equations (Spring 2019), Numerical Analysis
(Spring 2019), Nonlinear Differential Equations (Fall 2018), Numerical
Linear Algebra (Fall 2018), Mathematical Modeling with Stochastic
Differential Equations (Spring 2018)
\end{tabular}

%#################################################################################
\section*{Technical and Computational Skills}

\noindent
\begin{tabular}{lp{0.7\textwidth}}
Languages: & Python, C++, MATLAB, Julia \\
Numerical tools: & FEniCS, hIPPYlib, IPOPT, PyTorch, TensorFlow, JAX,
SciPy \\
Parallel programming: & OpenMP, MPI \\
Other: & Git, \LaTeX, Vim
\end{tabular}

\end{document}